%\setcounter{secnumdepth}{4} %Nummerieren bis in die 4. Ebene
%\setcounter{tocdepth}{4} %Inhaltsverzeichnis bis zur 4. Ebene

%\pagestyle{headings}


\sloppy % LaTeX ist dann nicht so streng mit der Silbentrennung
%~ \MakeShortVerb{\§}

\parindent0mm
\parskip0.5em


{
%\textwidth170mm
%\oddsidemargin30mm
%\evensidemargin30mm
%\addtolength{\oddsidemargin}{-1in}
%\addtolength{\evensidemargin}{-1in}

%\parskip0pt plus2pt

% Die Raender muessen eventuell fuer jeden Drucker individuell eingestellt
% werden. Dazu sind die Werte fuer die Abstaende `\oben' und `\links' zu
% aendern, die von mir auf jeweils 0mm eingestellt wurden.

%\newlength{\links} \setlength{\links}{10mm}  % hier abzuaendern
%\addtolength{\oddsidemargin}{\links}
%\addtolength{\evensidemargin}{\links}

\begin{titlepage}
\vspace*{-1.5cm}
\raisebox{16mm}{
    \begin{minipage}[t]{70mm}
        \begin{center}
            %\selectlanguage{german}
            {\Large AKAD Universität\\}
            {\normalsize
                AKAD Bildungsgesellschaft mbH\\
            }
            \vspace{3mm}
            {\small Heilbronner Straße 86 -- 70191 Stuttgart\\}
        \end{center}
    \end{minipage}
}
\hfill
\raisebox{0mm}{
    \includegraphics[width=150pt]{Images/Ikon.png}}
\vspace{10em}

% Titel
\begin{center}
%    \baselineskip=55pt
    %\baselineskip=45pt
    \textbf{\huge \titel}\\
    \ifthenelse{\isempty{\langtitel}}{}{
            {\LARGE \langtitel}
    }
    \baselineskip=0 pt
\end{center}

%\vspace{7em}

%\vfill

% Autor
%\begin{center}
%    \textbf{\Large
%        \bearbeiter
%    }
%\end{center}

\begin{center}
    \textbf{\Large
        \studiengang{}
    }
\end{center}

\vspace{15mm}

% Prüfungsordnungs-Angaben
\begin{center}
%\selectlanguage{german}

%%%%%%%%%%%%%%%%%%%%%%%%%%%%%%%%%%%%%%%%%%%%%%%%%%%%%%%%%%%%%%%%%%%%%%%%%
% Ja, richtig, hier kann die BA-Vorlage zur MA-Vorlage gemacht werden...
% (nicht mehr nötig!)
%%%%%%%%%%%%%%%%%%%%%%%%%%%%%%%%%%%%%%%%%%%%%%%%%%%%%%%%%%%%%%%%%%%%%%%%%
{\Large \arbeit}


\vspace{2em}
\ifthenelse{\equal{\sprache}{deutsch}}{
    \begin{tabular}[t]{ll}
    Vorgelegt von: & \bearbeiter\\
    Matrikelnummer: & \matrikelnummer\\
    E-Mail Adresse: & \email\\
    Modul: & \akadmodule\\
    Beginn des Assignments:& \beginndatum \\
    Abgabe des Assignments:& \abgabedatum \\
    Gutachter:         & \erstgutachter \\
    & \zweitgutachter \\
}{
    \begin{tabular}[t]{ll}
    Date of issue:& \beginndatum \\
    Date of submission:& \abgabedatum \\
    Reviewers:         & \erstgutachter \\
    & \zweitgutachter \\
}
\end{tabular}
\end{center}

\end{titlepage}

}


\clearpage
\begin{titlepage}
    \vspace*{\fill}

    \section*{Erklärung}


    Ich versichere, dass ich das vorliegende \arbeit{} selbstständig verfasst, keine
    anderen als die angegebenen Quellen und Hilfsmittel benutzt sowie alle wörtlich
    oder sinngemäß übernommenen Stellen in der Arbeit gekennzeichnet habe. Alle
    Teile meiner Arbeit, die wortwörtlich oder dem Sinn nach anderen Werken
    entnommen sind, wurden unter Angabe der Quelle Kenntnis gemacht. Gleiches gilt
    auch für Zeichnungen, Skizzen, bildliche Darstellung sowie für Quellen aus dem
    Internet, dazu zählen auch KI-basierte Anwendungen oder Werkzeuge. Die Arbeit
    wurde in gleicher oder ähnlicher Form noch nicht als Prüfungsleistung eingereicht.

    \vspace{25 mm}

    \begin{tabular}{lc}
        \ort{}, den \abgabedatum \hspace*{2cm} & \underline{\hspace{6cm}} \\
                                                   & \bearbeiter
    \end{tabular}

    \vspace*{\fill}
\end{titlepage}

%\ifthenelse{\equal{\zweiseitig}{twoside}}
\clearpage
\linespread{1}
\pagenumbering{Roman}
\tableofcontents
%\setcounter{page}{1}
%\pagenumbering{arabic}
\setcounter{tocdepth}{2}
